\chapter{Variables and Data Types}
\label{ch:variables-types}

In Chapter~1 we used a variable without dwelling on it. Now we slow down and look
properly at the two ideas that underpin every program: \emph{variables} (named
places to keep values) and \emph{types} (what kind of value a place may hold).
Get these right and everything else in the language falls into place.

\section{Variables: named boxes}

A \textbf{variable} is a name attached to a value. Picture a labelled box: the
label is the name, and the contents are the value. You create a variable by
writing its name, its type, and (usually) an initial value:

\begin{salam}
func main:
    age : int = 30
    name : str = "Sara"
    pi : f64 = 3.14159
    println name, "is", age, "years old"
end
\end{salam}

\begin{progout}
Sara is 30 years old
\end{progout}

Read each declaration left to right: a \emph{name}, then a \emph{type}, then
\code{=} and an \emph{initial value}. So \code{age int = 30} means ``make a
variable called \code{age}, of type \code{int} (a 32-bit integer), holding
\code{30}.''

\begin{note}
Unlike some languages, Salam does \emph{not} use a keyword such as \code{let} or
\code{var} to introduce a variable. The shape ``\textit{name} \textit{type}
\code{=} \textit{value}'' is enough. This keeps declarations short and lines up
the names neatly down the left edge of your code.
\end{note}

\section{Mutable and immutable variables}

By default, a variable in Salam is \textbf{immutable}: once you give it a value,
that value cannot change. This is a deliberate and very useful safety net it
means a value you thought was fixed really is fixed. Try to reassign an immutable
variable and the compiler stops you:

\begin{salam}
func main:
    score : int = 10
    score = 20        // error: 'score' is immutable
end
\end{salam}

When you genuinely need a value to change over time a counter, a running
total, a loop variable mark it \textbf{mutable} with the \code{mut} keyword:

\begin{salam}
func main:
    mut score : int = 10
    score = 20
    score = score + 5
    println score
end
\end{salam}

\begin{progout}
25
\end{progout}

There is also \code{const}, for a value that is fixed \emph{and} known when the
program is compiled a true constant, such as a maximum size or a
mathematical constant:

\begin{salam}
const MAX_USERS = 100

func main:
    println "capacity:", MAX_USERS
end
\end{salam}

\begin{progout}
capacity: 100
\end{progout}

So there are three flavours of binding, from most to least restrictive:

\begin{itemize}[leftmargin=1.4em]
  \item \code{const} fully immutable, and computed at compile time. Neither the
        binding nor its contents can change, not even a single array element or
        struct field. Use for true constants. Conventionally written in
        \code{UPPER\_CASE}.
  \item plain (\emph{bare}) the \emph{binding} is fixed once set: you cannot
        reassign the whole variable. But if it holds a compound value you may
        still change what is \emph{inside} it: array elements and struct fields
        stay writable. This is the default; prefer it.
  \item \code{mut} may be reassigned outright. Use only when a whole value
        really must be replaced.
\end{itemize}

The distinction between a bare binding and \code{const} is about depth. A bare
binding freezes the name; \code{const} freezes everything reachable through it:

\begin{salam}
func main:
    grid : int[3] = [1, 2, 3]
    grid[0] = 10          // fine: the binding is fixed, its elements are not
    // grid = [4, 5, 6]   // error: cannot reassign the whole binding

    const limits : int[3] = [1, 2, 3]
    // limits[0] = 10     // error: a const value is frozen through and through

    println grid[0]
end
\end{salam}

\begin{progout}
10
\end{progout}

\begin{tip}
Reach for \code{mut} only when you must. ``Immutable by default'' is one of
Salam's quiet strengths: the fewer things that can change, the fewer things can
go wrong, and the easier your program is to reason about. If the compiler
complains that something is immutable, pause and ask whether it truly needs to
change before adding \code{mut}.
\end{tip}

\section{Type inference with \texttt{auto}}

Often the type is obvious from the value, and writing it feels redundant. The
\code{auto} keyword asks the compiler to \emph{infer} the type from the
initializer:

\begin{salam}
func main:
    count : auto = 42          // inferred as int
    ratio : auto = 3.14        // inferred as f64
    label : auto = "items"     // inferred as str
    ready : auto = true        // inferred as bool
    println count, ratio, label, ready
end
\end{salam}

\begin{progout}
42 3.14 items true
\end{progout}

\code{auto} is about convenience, not weakness: the variable still has one fixed
type, decided once, at compile time. \code{count} above is exactly an \code{int}
--- you simply did not have to spell it out. You can combine \code{auto} with
\code{mut}: \code{mut total auto = 0} makes a mutable \code{int}.

\begin{warn}
\code{auto} needs an initializer to infer from. \code{x auto} on its own is
meaningless there is nothing to deduce the type from and the compiler
will reject it. When in doubt, or when you want to document intent, write the
type out explicitly.
\end{warn}

\section{The primitive types}

Salam's built-in scalar types are the atoms from which everything larger is
built. Here are the ones you will use constantly.

\subsection{Integers}

Integers are whole numbers. Salam offers them in several sizes, both
\emph{signed} (able to hold negatives) and \emph{unsigned} (zero and up only).
The number in the name is how many bits the value occupies, which fixes its
range:

\begin{center}
\begin{tabular}{@{}llll@{}}
\toprule
\textbf{Type} & \textbf{Signed?} & \textbf{Bits} & \textbf{Range} \\
\midrule
\code{i8}  & signed   & 8  & $-128$ \dots\ $127$ \\
\code{i16} & signed   & 16 & $-32768$ \dots\ $32767$ \\
\code{i32} & signed   & 32 & about $\pm 2.1$ billion \\
\code{i64} & signed   & 64 & about $\pm 9.2$ quintillion \\
\code{u8}  & unsigned & 8  & $0$ \dots\ $255$ \\
\code{u16} & unsigned & 16 & $0$ \dots\ $65535$ \\
\code{u32} & unsigned & 32 & $0$ \dots\ about $4.3$ billion \\
\code{u64} & unsigned & 64 & $0$ \dots\ about $18$ quintillion \\
\bottomrule
\end{tabular}
\end{center}

You do not need to memorise these. For everyday whole numbers, reach for
\code{int} it is the default and is almost always the right choice. Step up to
\code{i64} when you might exceed two billion (timestamps, file sizes, big
counters), and use the small or unsigned types only when you have a specific
reason, such as raw bytes (\code{u8}) or interfacing with C.

\begin{note}
\code{int} \emph{is} \code{i32}: a 32-bit signed integer. It is simply the
friendlier, everyday spelling of the same type, and the compiler treats the two
names as identical. The other two width-carrying types have short spellings as
well: \code{uint} is \code{u32}, and \code{float} is \code{f32}. This book writes
\code{int} and \code{float} in ordinary code, because the width is rarely the
point; it falls back on \code{i32}, \code{u32}, and \code{f32} in the places
where the number of bits genuinely matters, such as the table above or a C
interface. The compiler's own error messages always print the explicit names, so
do not be surprised to see \code{i32} looking back at you from a diagnostic.
\end{note}

\subsection{Floating-point numbers}

For numbers with a fractional part measurements, averages, scientific
values use a floating-point type:

\begin{itemize}[leftmargin=1.4em]
  \item \code{f64} a 64-bit ``double precision'' number. This is the default
        for decimal literals and the one to use unless you have a reason not to.
  \item \code{float} (the explicit spelling is \code{f32}) a 32-bit number;
        less precise, but smaller. Useful for graphics or large arrays where
        space matters.
\end{itemize}

\begin{salam}
func main:
    temperature : f64 = 36.6
    half : auto = 7.0 / 2.0        // 3.5, a floating-point division
    println temperature, half
end
\end{salam}

\begin{progout}
36.6 3.5
\end{progout}

\subsection{Booleans}

A \code{bool} holds one of just two values: \code{true} or \code{false}. Booleans
are the language of decisions every \code{if} and every \code{until} in the
coming chapters tests a \code{bool}.

\begin{salam}
func main:
    raining : bool = false
    weekend : auto = true
    println "raining?", raining
    println "weekend?", weekend
end
\end{salam}

\begin{progout}
raining? false
weekend? true
\end{progout}

\subsection{Characters}

A \code{char} is a single byte one character in the basic Latin range, or any
value from 0 to 255. Character literals are written between \emph{single} quotes:

\begin{salam}
func main:
    grade : char = 'A'
    println grade
    println grade as int        // the numeric code of 'A'
    println 66 as char          // the character with code 66
end
\end{salam}

\begin{progout}
A
65
B
\end{progout}

A \code{char} is really a small number in disguise: \code{'A'} \emph{is} 65. You
can convert between the two with \code{as}, which is how the example turns
\code{'A'} into \code{65} and \code{66} back into \code{'B'}.

\begin{note}
A single quote makes a \code{char}; a double quote makes a \code{str}. So
\code{'A'} is one character, while \code{"A"} is a string that happens to be one
character long. They are different types. We meet strings properly in
Chapter~8.
\end{note}

\subsection{Strings}

A \code{str} is a piece of text a sequence of characters written between
double quotes. Strings are \emph{immutable}: operations on them produce new
strings rather than changing the original. You can join two strings with
\code{+}:

\begin{salam}
func main:
    first : str = "Salam"
    greeting : auto = first + ", world!"
    println greeting
end
\end{salam}

\begin{progout}
Salam, world!
\end{progout}

\subsection{Void}

Finally, \code{void} is the ``no value'' type. You never store a \code{void};
it appears only as the result type of a function that returns nothing (a
function that \emph{does} something rather than \emph{computes} something). We
return to it in Chapter~6.

\section{Literals: writing values directly}

A \textbf{literal} is a value written straight into your program. Each kind of
literal has a natural type:

\begin{center}
\begin{tabular}{@{}lll@{}}
\toprule
\textbf{Literal} & \textbf{Type} & \textbf{Examples} \\
\midrule
whole number      & \code{int} & \code{0}, \code{42}, \code{-7} \\
decimal number    & \code{f64} & \code{3.14}, \code{2.0}, \code{1.5e3} \\
character          & \code{char} & \code{'a'}, \code{'9'} \\
text               & \code{str} & \code{"hi"}, \code{""} \\
truth value        & \code{bool} & \code{true}, \code{false} \\
\bottomrule
\end{tabular}
\end{center}

This leads to one rule that surprises newcomers and is worth learning now.

\begin{warn}
\textbf{A plain whole-number literal is an \code{int}.} Salam does not silently
shrink it to fit a smaller or unsigned type. So this is an \emph{error}:
\begin{salam}
mut b : u8 = 250           // error: cannot assign 'i32' to 'u8'
\end{salam}
The literal \code{250} is an \code{int}, and Salam will not narrow it for you
without being asked. (The compiler spells the type \code{i32} in its message;
that is the same type under its explicit name.) The fix is an explicit cast with
\code{as}:
\begin{salam}
mut b : u8 = 250 as u8     // fine
\end{salam}
This strictness is a feature: it makes every change of type visible in the
source, so a narrowing can never happen by accident.
\end{warn}

\section{Converting between types with \texttt{as}}

Salam never converts a value from one type to another behind your back. When you
\emph{do} want a conversion, you ask for it explicitly with the \code{as}
operator:

\begin{salam}
func main:
    total : int = 7
    count : int = 2
    average : f64 = (total as f64) / (count as f64)
    println "average =", average
end
\end{salam}

\begin{progout}
average = 3.5
\end{progout}

Why the casts? Because \code{total} and \code{count} are integers,
\code{total / count} would be \emph{integer} division it would compute
\code{3}, throwing away the remainder. By converting each to \code{f64} first, we
ask for real division and get \code{3.5}.

\begin{warn}
Integer division truncates toward zero: \code{7 / 2} is \code{3}, not
\code{3.5}, and \code{7 \% 2} (the remainder) is \code{1}. If you want a
fractional result, at least one operand must be a floating-point value. This is
one of the most common beginner surprises keep it in mind.
\end{warn}

\section{Type aliases}

When a type's meaning matters more than its mechanics, you can give it a clearer
name with \code{type}:

\begin{salam}
type Age = u32

func main:
    a : Age = 30 as u32
    println a
end
\end{salam}

\begin{progout}
30
\end{progout}

An alias is \emph{transparent}: \code{Age} and \code{u32} are completely
interchangeable the alias is just a more readable spelling. Aliases come into
their own with the bigger types we build later; for now, simply know the keyword
exists.

\section{Naming variables}

A good name is documentation that never goes out of date. Salam identifiers may
contain letters, digits, and underscores, and may not start with a digit.
Crucially, they may also contain \emph{Unicode} letters so Persian names are
first-class:

\begin{salam}
func main:
    user_count : int = 5
    maxScore : auto = 100
    println user_count, maxScore
end
\end{salam}

\begin{tip}
Choose names that say what the value \emph{means}, not what type it is.
\code{user\_count} tells the reader far more than \code{n} or \code{intval}. The
extra keystrokes pay for themselves the first time you re-read your own code a
week later.
\end{tip}

\section{Persian type names}

Just as the keywords have Persian forms, so do the types. A Persian program
writes the very same declarations with Persian type names:

\begin{center}
\begin{tabular}{@{}ll@{}}
\toprule
\textbf{English} & \textbf{Persian} \\
\midrule
\code{int} & \code{\fa{صحیح}} \\
\code{i64} & \code{\fa{صحیح۶۴}} \\
\code{f64} & \code{\fa{اعشار۶۴}} \\
\code{bool} & \code{\fa{منطقی}} \\
\code{char} & \code{\fa{کاراکتر}} \\
\code{str} & \code{\fa{رشته}} \\
\code{auto} & \code{\fa{خودکار}} \\
\bottomrule
\end{tabular}
\end{center}

\noindent Here is the chapter's first program written in Persian:

\begin{salamfa}
تابع اصلی:
    سن : صحیح = 30
    نام : رشته = "سارا"
    چاپ نام, "is", سن, "years old"
تمام
\end{salamfa}

Just as \code{int} is the everyday spelling in English, \code{\fa{صحیح}} is the
everyday spelling in Persian, and the two name exactly the same 32-bit type. When
you do want to state a width, the Persian name carries it in Persian digits:
\code{\fa{صحیح۶۴}} is \code{i64}, and the \fa{۶۴} inside it is part of the
keyword's spelling. Values, by contrast, are written with ordinary digits (the
\code{30} above), because arithmetic uses the standard digit characters. We will
untangle this fully in Chapter~22; for now, just admire that the structure is
identical to the English version.

\section{In Depth: how values are stored}

\begin{deep}[In Depth: representation and the generated C]
Each Salam scalar maps directly onto a C type when the program is compiled:
\code{i32} (that is, \code{int}) becomes \code{int32\_t}, \code{f64} becomes
\code{double},
\code{bool} becomes a one-byte value, \code{char} a single byte, and \code{str}
a pointer to immutable, NUL-terminated bytes. This 1-to-1 mapping is why Salam
programs run at C speed and why the foreign function interface (Chapter~19) is
so direct. Because the sizes are fixed and known, the compiler can check at
\emph{compile time} that \code{250} does not fit in an \code{i8} and that a
\code{u8} and an \code{i32} are not accidentally mixed the strictness you met
above is the visible edge of this static type system.
\end{deep}

\section{Chapter summary}

\begin{itemize}[leftmargin=1.4em]
  \item A variable is declared as \textit{name} \textit{type} \code{=}
        \textit{value}; no \code{let}/\code{var} keyword is needed.
  \item Variables are \emph{immutable by default}; add \code{mut} to allow
        reassignment, or \code{const} for a compile-time constant.
  \item \code{auto} infers the type from the initializer.
  \item Integers come in signed (\code{i8}--\code{i64}) and unsigned
        (\code{u8}--\code{u64}) sizes; the 32-bit one is the everyday default.
        \code{int}/\code{uint}/\code{float} are the short spellings of
        \code{i32}/\code{u32}/\code{f32}, and this book writes \code{int} and
        \code{float} in ordinary code.
  \item \code{f64} is the default for decimals; \code{bool} is \code{true} or
        \code{false}; \code{char} is one byte; \code{str} is immutable text.
  \item A whole-number literal is an \code{int}; conversions are never implicit ---
        use \code{as}. Integer division truncates.
  \item Types have Persian names too (\code{\fa{صحیح}}, \code{\fa{رشته}},
        \dots).
\end{itemize}

\section{Exercises}

\exercise Declare three variables describing a book: a title (\code{str}), a page
count (\code{int}), and a price (\code{f64}). Print them in one sentence.

\exercise Start with \code{mut total int = 0}, add \code{17} and then \code{25}
to it in two steps, and print the result after each step.

\exercise Compute the average of \code{10}, \code{15}, and \code{20} as an
\code{f64} and print it. (Remember what plain integer division would do.)

\exercise Store the character \code{'Z'} in a \code{char}, then print both the
character and its numeric code.

\exercise Try to assign \code{300} to a \code{u8} variable and read the compiler's
error. Then fix it with a cast and explain in a comment what the cast does.

\exercise Rewrite the book-description program from the first exercise using
Persian type names.
